\documentclass{zkdl-presentation-template}

% Circuit notation of the book
\providecommand{\Circ}{\textit{\texttt{C}}}

% Circuit diagrams (opencolor roles from CONVENTIONS.md)
\tikzset{
    var/.style={circle, draw=oc-gray-7, fill=oc-gray-2, very thick, minimum size=0.7cm, text centered},
    gate1/.style={circle, draw=oc-indigo-7, fill=oc-indigo-1, ultra thick, minimum size=1cm, text centered},
    gate2/.style={circle, draw=oc-orange-7, fill=oc-orange-1, ultra thick, minimum size=1cm, text centered},
    gate3/.style={circle, draw=oc-blue-7, fill=oc-blue-1, ultra thick, minimum size=1cm, text centered},
    arrow/.style={-latex, ultra thick},
    wire/.style={arrow, oc-gray-5},
    outwire/.style={arrow, oc-gray-7}
}

% PCP oracle diagrams
\tikzset{
    oracle/.style={fill=oc-gray-1, draw=oc-gray-9, ultra thick},
    query/.style={fill=oc-blue-4, draw=oc-gray-9, ultra thick}
}

% Plots
\pgfplotsset{
    zkdlgrid/.style={
        grid = both,
        grid style = {line width=.1pt, draw=oc-gray-2},
        major grid style = {line width=.2pt, draw=oc-gray-5}
    }
}

\title[QAP and PCP]{\textbf{QAP, PCP, PoE: Demystifying zk-SNARK Tools}}
\author{Blockstream Research}
\date{\today}
\homepage{research.blockstream.com}
\github{BlockstreamResearch/zk-book}

\begin{document}
    \maketitleframe

    \begin{frame}{Plan}
        \tableofcontents
    \end{frame}

    \section{Recap}

    \begin{frame}{Recap: what is zk-SNARK?}
        \begin{definition}
            \begin{center}
                \textbf{zk-SNARK} \\
                \textbf{Z}ero-\textbf{K}nowledge \textbf{S}uccinct \textbf{N}on-interactive \textbf{AR}gument of \textbf{K}nowledge.
            \end{center}
        \end{definition}
        \pause

        \begin{itemize}[label=\ding{51}]
            \item \textbf{Argument of Knowledge} --- a proof that the prover knows the data (witness) that resolves a certain
            problem, and this knowledge can be ``extracted''. \pause
            \item \textbf{Succinctness} --- the proof size and verification time is relatively small to the computation size and typically does not depend on the size of
            the data or statement. \pause
            \item \textbf{Non-interactiveness} --- to produce the proof, the prover does not need any interaction
            with the verifier. \pause
            \item \textbf{Zero-Knowledge} --- the verifier learns nothing about the data used to produce the
            proof, despite knowing that this data resolves the given problem and that the prover possesses it.
        \end{itemize}
    \end{frame}

    \begin{frame}{Recap: Arbitrary Program To Circuits}
        We can do that in a way like the computer does it --- \textbf{boolean circuits}.

        \begin{figure}[h!]
            \centering
            \begin{tabular}{cc}
                \begin{tikzpicture}[scale=0.85, every node/.style={scale=0.85}]
                    \node[var] (a) at (0, -1.5) {$a$};
                    \node[var] (b) at (2, -1.5) {$b$};
                    \node[gate1] (and) at (1, 0) {\texttt{AND}};
                    \node[var] (c) at (1, 1.75) {$c$};

                    \draw[wire] (a) -- (and);
                    \draw[wire] (b) -- (and);
                    \draw[outwire] (and) -- (c);
                \end{tikzpicture}
                &
                \begin{tikzpicture}[scale=0.85, every node/.style={scale=0.85}]
                    \node[var] (a) at (0, -1.5) {$a$};
                    \node[var] (b) at (2, -1.5) {$b$};
                    \node[gate2] (or) at (1, 0) {\texttt{OR}};
                    \node[var] (c) at (1, 1.75) {$c$};

                    \draw[wire] (a) -- (or);
                    \draw[wire] (b) -- (or);
                    \draw[outwire] (or) -- (c);
                \end{tikzpicture}
            \end{tabular}
            \caption{Boolean \texttt{AND} and \texttt{OR} Gates}
        \end{figure}
        \pause

        > \textbf{100000 gates} just for \texttt{SHA256}\ldots
        \pause

        But nothing stops us from using something more powerful instead of boolean values, gates.
    \end{frame}

    \begin{frame}{Recap: Arbitrary Program To Circuits}
        Similar to Boolean Circuits, the \textbf{Arithmetic Circuits} consist of gates and wires.
        \begin{itemize}
            \item \textbf{Wires}: elements of some finite field $\mathbb{F}$.
            \item \textbf{Gates}: field addition ($+$) and multiplication ($\times$).
        \end{itemize}
        \pause

        \begin{figure}[h!]
            \centering
            \begin{tabular}{cc}
                \begin{tikzpicture}
                    \node[var] (a) at (0, -1.5) {$a$};
                    \node[var] (b) at (2, -1.5) {$b$};
                    \node[gate1] (add) at (1, 0) {$+$};
                    \node[var] (c) at (1, 1.75) {$c$};

                    \draw[wire] (a) -- (add);
                    \draw[wire] (b) -- (add);
                    \draw[outwire] (add) -- (c);
                \end{tikzpicture}
                &
                \begin{tikzpicture}
                    \node[var] (a) at (0, -1.5) {$a$};
                    \node[var] (b) at (2, -1.5) {$b$};
                    \node[gate2] (mul) at (1, 0) {$\times$};
                    \node[var] (c) at (1, 1.75) {$c$};

                    \draw[wire] (a) -- (mul);
                    \draw[wire] (b) -- (mul);
                    \draw[outwire] (mul) -- (c);
                \end{tikzpicture}
            \end{tabular}
            \caption{Addition and Multiplication Gates}
        \end{figure}
    \end{frame}

    \begin{frame}[fragile]{Recap: Arbitrary Program To Circuits}
        \begin{example}
            How can we translate \texttt{if} statements?
            \begin{lstlisting}[language=Python,numbers=none]
def example(a: bool, b: F, c: F) -> F:
    if a:
        return b * c
    else:
        return b + c
            \end{lstlisting}
            \pause
            We can transform such a function into the next expression:
            \vspace{-8pt}
            \begin{equation*}
                r = a \times (b \times c) + (1 - a) \times (b + c)
                \vspace{-8pt}
            \end{equation*}
            \pause
            Corresponding equations for the circuit are:
            \vspace{-8pt}
            \begin{equation*}
                \begin{aligned}
                    r_1 &= b \times c, \quad &r_3 &= 1 - a, \quad &r_5 &= r_3 \times r_2 \\
                    r_2 &= b + c, \quad &r_4 &= a \times r_1, \quad &r &= r_4 + r_5
                \end{aligned}
            \end{equation*}
        \end{example}
    \end{frame}

    \begin{frame}{Recap: Arbitrary Program To Circuits}
         \begin{figure}[h!]
             \centering
             \begin{tikzpicture}
                 % Nodes
                 \node[var] (c) at (0.5, -3) {$c$};
                 \node[var] (b) at (0.5, -1.5) {$b$};
                 \node[var] (a) at (0.5, 0) {$a$};
                 \node[var] (one) at (0.5, 1.5) {$1$};

                 % b+c and b*c gates
                 \node[gate1] (b_plus_c) at (2.5, -1.5) {$+$};
                 \node[gate2] (b_times_c) at (2.5, -3.0) {$\times$};

                 \node[gate3] (one_minus_a) at (2.5, 0.75) {$-$};

                 % a*b*c and (1-a)(b+c) gates
                 \node[gate2] (a_times_b_times_c) at (5, -2.0) {$\times$};
                 \node[gate2] (one_minus_a_times_b_plus_c) at (5, -0.5) {$\times$};

                 % a*b*c + (1-a)(b+c) gate
                 \node[gate1] (r) at (7.5, -1.25) {$+$};

                 % Result node
                 \node[var] (result) at (10, -1.25) {$r$};

                 % b+c and b*c arrows
                 \draw[wire] (b) to (b_plus_c);
                 \draw[wire] (b) to (b_times_c);
                 \draw[wire] (c) to (b_plus_c);
                 \draw[wire] (c) to (b_times_c);

                 % 1 - a arrows
                 \draw[wire] (one) to (one_minus_a);
                 \draw[wire] (a) to (one_minus_a);

                 % a*b*c and (1-a)(b+c) arrows
                 \draw[wire] (a) to [bend left=20] (a_times_b_times_c);
                 \draw[wire] (b_times_c) to node[midway, above] {$r_1$} (a_times_b_times_c);
                 \draw[wire] (one_minus_a) to node[midway, above] {$r_3$} (one_minus_a_times_b_plus_c);
                 \draw[wire] (b_plus_c) to node[midway, above] {$r_2$} (one_minus_a_times_b_plus_c);

                 % a*b*c + (1-a)(b+c) arrows
                 \draw[wire] (a_times_b_times_c) to [bend right=20] node[midway, above] {$r_4$} (r);
                 \draw[wire] (one_minus_a_times_b_plus_c) to [bend left=20] node[midway, above] {$r_5$} (r);

                 % Result arrow
                 \draw[outwire] (r) to (result);
             \end{tikzpicture}
             \caption{Example of a circuit evaluating the \texttt{if} statement logic.}
         \end{figure}
    \end{frame}

    \begin{frame}{Recap: R1CS}
        Each \textbf{constraint} in the Rank-1 Constraint System must be in the form:
        \begin{equation*}
            \langle \boldsymbol{\ell}, \mathbf{z}\rangle \times \langle \boldsymbol{r}, \mathbf{z}\rangle = \langle \boldsymbol{o}, \mathbf{z}\rangle
        \end{equation*}
        \pause
        Where $\langle \mathbf{u}, \mathbf{v}\rangle$ is a dot product.
        \vspace{-10pt}
        \begin{equation*}
            \langle \mathbf{u}, \mathbf{v} \rangle := \mathbf{u}^{\top}\mathbf{v} = \sum_{i=1}^{n} u_i v_i
            \vspace{-5pt}
        \end{equation*}
        \pause
        Thus
        \vspace{-5pt}
        \begin{equation*}
            \left(\sum_{i=1}^{n} \ell_i z_i\right) \times \left(\sum_{j=1}^{n} r_j z_j\right) = \sum_{k=1}^{n} o_k z_k
        \end{equation*}
        That is, actually, a quadratic equation with multiple variables.
    \end{frame}

    \begin{frame}{Recap: R1CS}
        \begin{example}
            Consider the most basic circuit with one multiplication gate: $x_1 \times x_2 = r$.
            The witness vector $\mathbf{z} = (r, x_1, x_2)$. So
            \vspace{-5pt}
            \begin{align*}
                z_2 &\times z_3 = z_1 \\
                (0 + z_2 + 0) &\times (0 + 0 + z_3) = z_1 + 0 + 0 \\
                (0z_1 + 1z_2 + 0z_3) &\times (0z_1 + 0z_2 + 1z_3) = 1z_1 + 0z_2 + 0z_3
            \end{align*}
            \pause
            Therefore the coefficients vectors are:
            \vspace{-5pt}
            \begin{equation*}
                \boldsymbol{\ell} = (0, 1, 0), \quad \boldsymbol{r} = (0, 0, 1), \quad \boldsymbol{o} = (1, 0, 0).
                \vspace{-5pt}
            \end{equation*}
            \pause
            The general form of our constraint is:
            \vspace{-5pt}
            \begin{equation*}
                (\ell_1z_1 + \ell_2z_2 + \ell_3z_3)(r_1z_1 + r_2z_2 + r_3z_3) = o_1z_1 + o_2z_2 + o_3z_3
            \end{equation*}
        \end{example}
    \end{frame}

    \begin{frame}{Recap: R1CS}
        \vspace{-10pt}
        \begin{equation*}
            r = x_1 \times (x_2 \times x_3) + (1 - x_1) \times (x_2 + x_3)
        \end{equation*}
        \pause
        Thus, the next constraints can be built:
        \vspace{-5pt}
        \begin{align*}
            x_1 \times x_1 &= x_1 \quad \text{(binary check)} \tag{1} \\
            x_2 \times x_3 &= \mathsf{mult} \tag{2} \\
            x_1 \times \mathsf{mult} &= \mathsf{selectMult} \tag{3} \\
            (1 - x_1) \times (x_2 + x_3) &= r - \mathsf{selectMult} \tag{4}
        \end{align*}
        \pause
        The witness vector: $\mathbf{z} = (1, r, x_1, x_2, x_3, \mathsf{mult}, \mathsf{selectMult})$.
        \pause

        \vspace{2pt}
        The coefficients vectors:
        \vspace{-25pt}
        {\center\small\begin{align*}
            \boldsymbol{\ell}_1 &= (0, 0, 1, 0, 0, 0, 0), & \boldsymbol{r}_1 &= (0, 0, 1, 0, 0, 0, 0), & \boldsymbol{o}_1 &= (0, 0, 1, 0, 0, 0, 0) \\
            \boldsymbol{\ell}_2 &= (0, 0, 0, 1, 0, 0, 0), & \boldsymbol{r}_2 &= (0, 0, 0, 0, 1, 0, 0), & \boldsymbol{o}_2 &= (0, 0, 0, 0, 0, 1, 0) \\
            \boldsymbol{\ell}_3 &= (0, 0, 1, 0, 0, 0, 0), & \boldsymbol{r}_3 &= (0, 0, 0, 0, 0, 1, 0), & \boldsymbol{o}_3 &= (0, 0, 0, 0, 0, 0, 1) \\
            \boldsymbol{\ell}_4 &= (1, 0, -1, 0, 0, 0, 0), & \boldsymbol{r}_4 &= (0, 0, 0, 1, 1, 0, 0), & \boldsymbol{o}_4 &= (0, 1, 0, 0, 0, 0, -1)
        \end{align*}}
    \end{frame}

    \section[QAP]{Quadratic Arithmetic Program}

    \begin{frame}{Why do we need QAP?}
        Problems we have for now: \pause
        \begin{itemize}[label=\ding{51}]
            \item Although Rank-1 Constraint Systems provide a powerful method for representing
            computations, they are not succinct. \pause
            \item We need to transform our computations into a form that is more convenient for
            proving statements about them.
        \end{itemize}
        \pause

        \begin{alertblock}{Notice}
            A very convenient form for representing computations is \textbf{polynomials}!

            \vspace{10px}

            \textbf{Idea:} Instead of checking polynomial equality $P(T)=Q(T)$ at multiple points $x_1,\dots,x_n$ (essentially, checking each constraint), we check it only once at $\tau \leftarrowS \mathbb{F}$: $P(\tau)=Q(\tau)$. Soundness is guaranteed by the \textbf{Schwartz-Zippel Lemma}.
        \end{alertblock}
    \end{frame}

    \begin{frame}{R1CS in Matrix Form}
        We finished with the following constraint vectors:
        \begin{equation*}
            \boldsymbol{\ell}_1, \boldsymbol{\ell}_2, \dots, \boldsymbol{\ell}_m, \quad
            \boldsymbol{r}_1, \boldsymbol{r}_2, \dots, \boldsymbol{r}_m, \quad
            \boldsymbol{o}_1, \boldsymbol{o}_2, \dots, \boldsymbol{o}_m.
        \end{equation*}
        \pause
        Of course, they form corresponding matrices:
        \begin{equation*}
            L = \begin{bmatrix}
                L_{1,1} & L_{1,2} & \dots & L_{1,n} \\
                L_{2,1} & L_{2,2} & \dots & L_{2,n} \\
                \vdots & \vdots & \ddots & \vdots \\
                L_{m,1} & L_{m,2} & \dots & L_{m,n}
            \end{bmatrix}, \;
            \text{same goes for $R$ and $O$}
        \end{equation*}
        \pause
        An example of a single ``if'' statement:
        \begin{center}
        \begin{minipage}{0.4\textwidth}
        \vspace{-15pt}
        {\small \begin{align*}
            \boldsymbol{\ell}_1 &= (0, 0, 1, 0, 0, 0, 0) \\
            \boldsymbol{\ell}_2 &= (0, 0, 0, 1, 0, 0, 0) \\
            \boldsymbol{\ell}_3 &= (0, 0, 1, 0, 0, 0, 0) \\
            \boldsymbol{\ell}_4 &= (1, 0, -1, 0, 0, 0, 0)
        \end{align*}}
        \end{minipage}
        \begin{minipage}{0.5\textwidth}
        \begin{equation*}
            L = {\small \begin{bmatrix}
                0 & 0 & 1 & 0 & 0 & 0 & 0 \\
                0 & 0 & 0 & 1 & 0 & 0 & 0 \\
                0 & 0 & 1 & 0 & 0 & 0 & 0 \\
                1 & 0 & -1 & 0 & 0 & 0 & 0
            \end{bmatrix}}
        \end{equation*}
        \end{minipage}
        \end{center}
        \pause

        Pleeeeeenty of zeroes, right? And this is just one out of 3 matrices...
    \end{frame}

    \begin{frame}{Columns of the Matrix}
        The previous witness vector:
        \begin{equation*}
            \mathbf{z} = (1, r, \overset{\textcolor{oc-red-7}{3}}{\textcolor{oc-red-7}{x_1}}, x_2, x_3, \mathsf{mult}, \mathsf{selectMult})
        \end{equation*}
        Let us take a closer look at the matrix columns:
        \begin{equation*}
            \begin{bmatrix}
                0 & 0 & \textcolor{oc-red-7}{1} & 0 & 0 & 0 & 0 \\
                0 & 0 & \textcolor{oc-red-7}{0} & 1 & 0 & 0 & 0 \\
                0 & 0 & \textcolor{oc-red-7}{1} & 0 & 0 & 0 & 0 \\
                1 & 0 & \textcolor{oc-red-7}{-1} & 0 & 0 & 0 & 0
            \end{bmatrix}
        \end{equation*}
        \pause
        Consider 4th constraint: $(1 - x_1) \times (x_2 + x_3) = r - \mathsf{selectMult}$
        \begin{equation*}
            \begin{bmatrix}
                0 & 0 & \textcolor{oc-red-7}{1} & 0 & 0 & 0 & 0 \\
                0 & 0 & \textcolor{oc-red-7}{0} & 1 & 0 & 0 & 0 \\
                0 & 0 & \textcolor{oc-red-7}{1} & 0 & 0 & 0 & 0 \\
                \textcolor{oc-blue-7}{1} & \textcolor{oc-blue-7}{0} & \boxed{-1} & \textcolor{oc-blue-7}{0} & \textcolor{oc-blue-7}{0} & \textcolor{oc-blue-7}{0} & \textcolor{oc-blue-7}{0}
            \end{bmatrix}
        \end{equation*}
        \pause
        So, every column is a mapping of constraint number to a coefficient for the witness element.
    \end{frame}

    \begin{frame}{Polynomial Interpolation}
        As we know, such a mapping can be built using Lagrange interpolation polynomial with the
        following formula:
        \vspace{-10pt}
        \begin{equation*}
            f(T) = \sum_{i=0}^{n} y_i e_i(T), \quad e_i(T) = \prod_{j=0, j \neq i}^{n} \frac{T-x_j}{x_i-x_j}.
            \vspace{-10pt}
        \end{equation*}
        \pause

        There are $n$ columns and $m$ constraints. So, it results in $n$ polynomials such that:
        \vspace{-10pt}
        \begin{equation*}
            \ell_j(i) = L_{i,j}, \; i \in \{1,\dots,m\}, \; j \in \{1,\dots,n\}
            \vspace{-10pt}
        \end{equation*}
        \pause

        The same is true for matrices $R$ and $O$, with $3n$ polynomials in total, $n$ for each of the
        coefficients matrices:
        \vspace{-8pt}
        \begin{equation*}
            \ell_1(T), \dots, \ell_n(T),
            r_1(T), \dots, r_n(T),
            o_1(T), \dots, o_n(T)
        \end{equation*}
        \pause

        \vspace{-10pt}
        \begin{remarkblock}{Note}
            We could have assigned any \textit{unique} index from $\mathbb{F}$ to each constraint
            (say, $t_i$ for each $i \in \{1,\dots,m\}$) and interpolate through these points:
            \vspace{-8pt}
            \begin{equation*}
                \ell_j(t_i) = L_{i,j}, \; i \in \{1,\dots,m\}, \; j \in \{1,\dots,n\}
            \end{equation*}
        \end{remarkblock}
    \end{frame}

    \begin{frame}{Polynomial Interpolation: Example}
        \begin{example}
            Considering the witness vector $\mathbf{z}$ and matrix $L$ from the previous example, for the variable
            $x_1$, the next set of points can be derived:
            \vspace{-10pt}
            \begin{equation*}
                \{(1,1), (2,0), (3,1), (4,-1)\}
            \end{equation*}
            \pause
            The Lagrange interpolation polynomial for this set of points:
            \begin{align*}
                e_1(T) &= -\frac{(T - 2)(T - 3)(T - 4)}{6}, \; &e_2(T) &= \frac{(T - 1)(T - 3)(T - 4)}{2}, \\
                e_3(T) &= -\frac{(T - 1)(T - 2)(T - 4)}{2}, \; &e_4(T) &= \frac{(T - 1)(T - 2)(T - 3)}{6}.
            \end{align*}
            \pause

            Thus, the polynomial is given by:
            \begin{align*}
                \ell_{x_1}(T) &= 1 \cdot e_1(T) + 0 \cdot e_2(T) + 1 \cdot e_3(T) + (-1) \cdot e_4(T) \\
                &= -\frac{5}{6} T^3 + 6T^2 - \frac{79}{6} T + 9
            \end{align*}
        \end{example}
    \end{frame}

    \begin{frame}{Polynomial Interpolation: Example}
        \begin{center}
            \begin{tikzpicture}
                \begin{axis}[
                    axis lines = middle,
                    xlabel = {$T$},
                    ylabel = {$\ell_{x_1}(T)$},
                    ymin = -2.99, ymax = 2.99,
                    xmin = -0.5, xmax = 4.99,
                    domain = 0:5,
                    samples = 100,
                    ytick = {-3,...,3},
                    xtick = {0,1,...,5},
                    zkdlgrid
                ]
                \addplot[
                    color=oc-blue-7,
                    thick
                ]
                {-5/6*x^3 + 6*x^2 - 79/6*x + 9};
                \addplot[
                    only marks,
                    mark=*,
                    color=oc-red-7
                ]
                coordinates {(1,1) (2,0) (3,1) (4,-1)};

                \node at (axis cs:1,1.5) [anchor=west] {\small $(1,1)$};
                \node at (axis cs:2,0.2) [anchor=south] {\small $(2,0)$};
                \node at (axis cs:3,1.5) [anchor=west] {\small $(3,1)$};
                \node at (axis cs:4,-1) [anchor=west] {\small $(4,-1)$};

                \end{axis}
            \end{tikzpicture}

            \scriptsize\textbf{Illustration:} The Lagrange interpolation polynomial for points $\{(1,1), (2,0), (3,1), (4,-1)\}$ visualized over $\mathbb{R}$.
        \end{center}
    \end{frame}

    \begin{frame}{Why Polynomials?}
        \begin{questionblock}{Question}
            But what does it change? We ``exchanged'' $3n$ columns for $3n$ polynomials.
        \end{questionblock}
        \pause

        \vspace{-5pt}
        Consider two polynomials $p(T)$ and $q(T)$:
        \begin{equation*}
            \quad p(T) = -\frac{1}{2}T^2 + \frac{3}{2}T, \qquad
            q(T) = \frac{1}{3}T^3 - 2T^2 + \frac{8}{3}T + 1.
            \vspace{-5pt}
        \end{equation*}
        With corresponding sets of points:
        \vspace{-5pt}
        \begin{equation*}
            \{(0, 0), (1, 1), (2, 1), (3, 0)\}, \quad \{(0, 1), (1, 2), (2, 1), (3, 0)\}
            \vspace{-10pt}
        \end{equation*}
        \pause

        The sum of these polynomials can be calculated as:
        \vspace{-5pt}
        \begin{equation*}
            (p+q)(T) = \frac{1}{3}T^3 - \frac{5}{2}T^2 + \frac{25}{6}T + 1
            \vspace{-5pt}
        \end{equation*}
        The resulting polynomial $p+q$ corresponds to the set of points:
        \vspace{-5pt}
        \begin{equation*}
            \{(0, 1), (1, 3), (2, 2), (3, 0)\}
        \end{equation*}
    \end{frame}

    \begin{frame}{Why Polynomials?}
        \begin{figure}[H]
            \centering
            \begin{tikzpicture}
                \begin{axis}[
                    axis lines = middle,
                    xlabel = {$T$},
                    ylabel = {$y$},
                    ymin = -0.99, ymax = 3.99,
                    xmin = -0.5, xmax = 3.99,
                    domain = 0:5,
                    samples = 100,
                    ytick = {-1,...,4},
                    xtick = {0,1,...,4},
                    zkdlgrid
                ]

                % p(T)
                \addplot[
                    color=oc-blue-7,
                    thick
                ]
                {3/2*x - 1/2*x^2};
                \addplot[
                    only marks,
                    mark=*,
                    color=oc-red-7
                ]
                coordinates {(0, 0) (1, 1) (2, 1) (3, 0)};
                \node at (axis cs:3.35,-0.7) [anchor=east] {$p(T)$};

                \node at (axis cs:0,-0.2) [anchor=west] {\tiny $(0, 0)$};
                \node at (axis cs:1,1) [anchor=south] {\tiny $(1, 1)$};
                \node at (axis cs:2.1,1) [anchor=south] {\tiny $(2, 1)$};
                \node at (axis cs:2.9,0.2) [anchor=west] {\tiny $(3, 0)$};

                % q(T)
                \addplot[
                    color=oc-green-7,
                    thick
                ]
                {1/3*x^3 - 2*x^2 + 8/3*x + 1};
                \addplot[
                    only marks,
                    mark=*,
                    color=oc-red-7
                ]
                coordinates {(0,1) (1,2) (2,1) (3, 0)};
                \node at (axis cs:3.7,0.7) [anchor=south] {$q(T)$};

                \node at (axis cs:0,1) [anchor=west] {\tiny $(0, 1)$};
                \node at (axis cs:1,2) [anchor=south] {\tiny $(1, 2)$};

                % (p+q)(T)
                \addplot[
                    color=oc-orange-7,
                    thick
                ]
                {1/3*x^3 - 5/2*x^2 + 25/6*x + 1};
                \addplot[
                    only marks,
                    mark=*,
                    color=oc-red-7
                ]
                coordinates {(0,1) (1,3) (2,2) (3, 0)};
                \node at (axis cs:3.4,-0.55) [anchor=west] {$(p+q)(T)$};

                \node at (axis cs:1,3) [anchor=south] {\tiny $(1, 3)$};
                \node at (axis cs:2,2) [anchor=west] {\tiny $(2, 2)$};

                \end{axis}
            \end{tikzpicture}
            \caption{Addition of two polynomials}
        \end{figure}
    \end{frame}

    \begin{frame}{Putting It All Together}
        \vspace{-5pt}
        Now, using coefficients encoded with polynomials, we can build a constraint number
        $j \in \{1, \dots, m\}$ in the next way:
        \vspace{-5pt}
        \begin{align*}
            &(z_1\ell_1(j) + z_2\ell_2(j) + \dots + z_n\ell_n(j)) \times \\
            \times &(z_1r_1(j) + z_2r_2(j) + \dots + z_nr_n(j)) = \\
            = &(z_1o_1(j) + z_2o_2(j) + \dots + z_no_n(j))
            \vspace{-5pt}
        \end{align*}
        \pause
        Or written more concisely:
        \vspace{-5pt}
        \begin{equation*}
            \left( \sum_{i = 1}^{n} z_i\ell_i(T) \right) \times \left( \sum_{i = 1}^{n} z_ir_i(T) \right) = \left( \sum_{i = 1}^{n} z_io_i(T) \right)
        \end{equation*}
    \end{frame}

    \begin{frame}{Putting It All Together}
        Hold on, but why does it hold? Let us substitute any $T=j$ into this equation:
        \vspace{-5pt}
        \begin{equation*}
            \left( \sum_{i = 1}^{n} z_i\ell_i(j) \right) \times \left( \sum_{i = 1}^{n} z_ir_i(j) \right) = \left( \sum_{i = 1}^{n} z_io_i(j) \right) \; \commentgray{for all $j \in \{1,\dots,m\}$}
            \vspace{-5pt}
        \end{equation*}
        \pause
        Recall that we interpolated polynomials to have $\ell_i(j) = L_{j,i}$. Therefore, the equation above can be reduced to:
        \vspace{-5pt}
        \begin{equation*}
            \left( \sum_{i = 1}^{n} z_iL_{j,i} \right) \times \left( \sum_{i = 1}^{n} z_iR_{j,i} \right) = \left( \sum_{i = 1}^{n} z_iO_{j,i} \right) \; \commentgray{for all $j \in \{1,\dots,m\}$}
            \vspace{-5pt}
        \end{equation*}
        \pause
        But hold on again! Notice that $\sum_{i = 1}^{n} z_iL_{j,i} = \langle \mathbf{z}, \boldsymbol{\ell}_j \rangle$ and therefore we have:
        \vspace{-5pt}
        \begin{equation*}
            \langle \mathbf{z}, \boldsymbol{\ell}_j \rangle \times \langle \mathbf{z}, \boldsymbol{r}_j \rangle = \langle \mathbf{z}, \boldsymbol{o}_j \rangle \; \commentgray{for all $j \in \{1,\dots,m\}$,}
            \vspace{-5pt}
        \end{equation*}

        so we ended up with the initial $m$ constraint equations!
    \end{frame}

    \begin{frame}{Vanishing Polynomial}
        Now let us define polynomials $\ell(T)$, $r(T)$, $o(T)$ for easier notation:
        \vspace{-5pt}
        \begin{equation*}
            \ell(T) = \sum_{i = 1}^{n} z_i\ell_i(T), \quad r(T) = \sum_{i = 1}^{n} z_ir_i(T), \quad o(T) = \sum_{i = 1}^{n} z_io_i(T)
            \vspace{-5pt}
        \end{equation*}
        \pause
        Therefore:
        \vspace{-5pt}
        \begin{equation*}
            \ell(j) \times r(j) = o(j) \; \commentgray{for all $j \in \{1,\dots,m\}$}
            \vspace{-5pt}
        \end{equation*}

        Now, we can define a polynomial $m(T)$, that has zeros at all elements from the set
        $\Omega = \{1,\dots,m\}$
        \vspace{-5pt}
        \begin{equation*}
            m(T) = \ell(T) \times r(T) - o(T)
            \vspace{-5pt}
        \end{equation*}
        \pause
        It means, that $m(T)$ can be divided by \textbf{vanishing polynomial} $z_{\Omega}(T)$ without a remainder!
        \vspace{-8pt}
        \begin{equation*}
            z_{\Omega}(T) = \prod_{i=1}^m (T - i), \qquad \text{$h(T) = \frac{m(T)}{z_{\Omega}(T)}$ is a polynomial}
            \vspace{-5pt}
        \end{equation*}
    \end{frame}

    \begin{frame}{Quadratic Arithmetic Program}
        \begin{definition}[Quadratic Arithmetic Program]
            Suppose that $m$ R1CS constraints with a witness of size $n$ are written in a form $L\mathbf{z} \odot R\mathbf{z} = O\mathbf{z}$ for $L,R,O \in \mathbb{F}^{m \times n}$.

            Then, the \textbf{Quadratic Arithmetic Program} consists of $3n$ polynomials $\ell_1,\dots,\ell_n$, $r_1,\dots,r_n$, $o_1,\dots,o_n$ such that:
            \begin{equation*}
                \ell_j(i) = L_{i,j}, \; r_j(i) = R_{i,j}, \; o_j(i) = O_{i,j} \; \commentgray{for all $i \in \{1,\dots,m\}$, $j \in \{1,\dots,n\}$}
            \end{equation*}

            Then, $\mathbf{z} \in \mathbb{F}^n$ is a valid assignment for the given QAP and \textbf{target polynomial} $z(T) = \prod_{i=1}^m (T-i)$ if and only if there exists such a polynomial $h(T)$ such that
            \begin{equation*}
                \left( \sum_{i = 1}^{n} z_i\ell_i(T) \right)\left( \sum_{i = 1}^{n} z_ir_i(T) \right) - \left( \sum_{i = 1}^{n} z_io_i(T) \right) = z(T)h(T)
            \end{equation*}
        \end{definition}
    \end{frame}

    \section{Probabilistically Checkable Proofs}

    \begin{frame}{Probabilistically Checkable Proofs}
        \begin{figure}[H]
            \centering
            \begin{tikzpicture}
                \node[inner sep=0pt, align=center] at (-4.0, 0.0) (prover) {\includegraphics[width=1.25cm]{images/common/prover.png}\\Prover $\mathcal{P}$};
                \node[inner sep=0pt, align=center] at (4.0, 0.0) (verifier) {\includegraphics[width=1.25cm]{images/common/verifier.png}\\Verifier $\mathcal{V}$};

                \draw[oracle] (-2.0, 4) rectangle (2.0, 3.25);

                \node at (0, 2.85) (pcp) {\textbf{PCP Oracle}};
                \draw[-{Stealth[length=3mm]},line width=0.4mm] (prover)--(pcp) node[midway, left=3mm]{\shortstack{Generate an\\oracle ($\pi$)}};
                \draw[-{Stealth[length=3mm]},line width=0.4mm] (verifier)--(pcp) node[midway, right=3mm]{\shortstack{Point queries\\$q_1,\dots,q_m$}};

                \pause
                % Query boxes
                \draw[query] (-1.0, 4) rectangle (-0.8, 3.25);
                \node at (-0.9, 4.35) {$q_1$};
                \draw[query] (0.0, 4) rectangle (0.2, 3.25);
                \node at (0.1, 4.35) {$q_2$};
                \draw[query] (0.5, 4) rectangle (0.7, 3.25);
                \node at (0.6, 4.35) {$q_3$};
            \end{tikzpicture}
            \caption{Illustration of a Probabilistically Checkable Proof (PCP) system. The prover $\mathcal{P}$ generates a PCP oracle $\pi$ that is queried by the verifier $\mathcal{V}$ at specific points $q_1,\dots,q_m$.}
        \end{figure}
    \end{frame}

    \begin{frame}{Extensions of PCPs}
        Three main extensions of PCPs that are frequently used in SNARKs are:
        \begin{itemize}
            \item \textbf{IPCP} (\textbf{Interactive PCP}): The prover commits to the PCP oracle and then, based on the interaction between the prover and verifier, the verifier queries the oracle and decides whether to accept the proof.\pause
            \item \textbf{IOP} (\textbf{Interactive Oracle Proof}): The prover and verifier interact and on each round, the prover commits to a new oracle. The verifier queries the oracle and decides whether to accept the proof.\pause
            \item \textbf{LPCP} (\textbf{Linear PCP}): The prover commits to a linear function and the verifier queries the function at specific points.
        \end{itemize}
    \end{frame}

    \begin{frame}{Interactive Oracle Proofs}
        \begin{figure}[H]
            \centering
            \begin{tikzpicture}[scale=0.9, every node/.style={scale=0.9}]
                \node[inner sep=0pt, align=center] at (-4.0, 0.5) (prover) {\includegraphics[width=1cm]{images/common/prover.png}\\Prover $\mathcal{P}$};
                \node[inner sep=0pt, align=center] at (4.0, 0.5) (verifier) {\includegraphics[width=1cm]{images/common/verifier.png}\\Verifier $\mathcal{V}$};

                % Oracle 1
                \node at (-2.0, 5.625) (oracle_1_left) {};
                \node at (2.0, 5.625) (oracle_1_right) {};
                \draw[oracle] (-2.0, 6) rectangle (2.0, 5.25);
                \draw[query] (-1.0, 6) rectangle (-0.8, 5.25);
                \node at (-0.9, 6.35) {$q_1$};
                \draw[query] (0.0, 6) rectangle (0.2, 5.25);
                \node at (0.1, 6.35) {$q_2$};
                \draw[query] (0.5, 6) rectangle (0.7, 5.25);
                \node at (0.6, 6.35) {$q_3$};
                \node at (0, 4.85) {\textbf{PCP Oracle \#1}};

                % Oracle 2
                \node at (-2.0, 3.625) (oracle_2_left) {};
                \node at (2.0, 3.625) (oracle_2_right) {};
                \draw[oracle] (-2.0, 4) rectangle (2.0, 3.25);
                \draw[query] (-0.4, 4) rectangle (-0.2, 3.25);
                \node at (-0.3, 4.35) {$q_1'$};
                \draw[query] (0.2, 4) rectangle (0.4, 3.25);
                \node at (0.3, 4.35) {$q_2'$};
                \draw[query] (1.5, 4) rectangle (1.7, 3.25);
                \node at (1.6, 4.35) {$q_3'$};
                \node at (0, 2.85) {\textbf{PCP Oracle \#2}};

                % Oracle 3
                \node at (-2.0, 1.625) (oracle_3_left) {};
                \node at (2.0, 1.625) (oracle_3_right) {};
                \draw[oracle] (-2.0, 2) rectangle (2.0, 1.25);
                \draw[query] (-1.5, 2) rectangle (-1.3, 1.25);
                \node at (-1.4, 2.35) {$q_1''$};
                \draw[query] (0.3, 2) rectangle (0.5, 1.25);
                \node at (0.4, 2.35) {$q_2''$};
                \draw[query] (1.0, 2) rectangle (1.2, 1.25);
                \node at (1.1, 2.35) {$q_3''$};
                \node at (0, 0.85) {\textbf{PCP Oracle \#3}};

                \draw[-{Stealth[length=3mm]},line width=0.4mm] (prover)--(oracle_1_left) node[midway, left=3mm]{\shortstack{Commit to\\oracles}};
                \draw[-{Stealth[length=3mm]},line width=0.4mm] (verifier)--(oracle_1_right) node[midway, right=3mm]{Point queries};
                \draw[-{Stealth[length=3mm]},line width=0.4mm] (prover)--(oracle_2_left);
                \draw[-{Stealth[length=3mm]},line width=0.4mm] (verifier)--(oracle_2_right);
                \draw[-{Stealth[length=3mm]},line width=0.4mm] (prover)--(oracle_3_left);
                \draw[-{Stealth[length=3mm]},line width=0.4mm] (verifier)--(oracle_3_right);

                % Two arrows to left and to right between the prover and verifier with the "Interaction" label
                \draw[{Stealth[length=3mm]}-{Stealth[length=3mm]},line width=0.4mm,dashed] (prover) to node[midway, below]{Interaction} (verifier);
            \end{tikzpicture}
            \caption{Illustration of an Interactive Oracle Proof (IOP). On each round $i$ ($1 \leq i \leq r$), $\mathcal{V}$ sends a message $m_i$, and $\mathcal{P}$ commits to a new oracle $\pi_i$, which $\mathcal{V}$ can query at $\mathbf{q}_i=(q_{i,1},\dots,q_{i,m})$.}
        \end{figure}
    \end{frame}

    \section{QAP as a Linear PCP}

    \begin{frame}{Linear PCP}
        \begin{definition}[Linear PCP]
            A \textbf{Linear PCP} is a PCP where the prover commits to a linear function $\boldsymbol{\pi} = (\pi_1,\dots,\pi_k)$ and the verifier queries the function at specific points $\mathbf{q}_1,\dots,\mathbf{q}_r$. Then, the prover responds with the values of the function at these points:
            \begin{equation*}
                \langle \boldsymbol{\pi}_1, \mathbf{q}_1 \rangle, \langle \boldsymbol{\pi}_2, \mathbf{q}_2 \rangle, \dots, \langle \boldsymbol{\pi}_r, \mathbf{q}_r \rangle.
            \end{equation*}
        \end{definition}
    \end{frame}

    \begin{frame}{QAP as a Linear PCP}
        \begin{example}[QAP as a Linear PCP]
            Recall that key QAP equation is:
            \begin{equation*}
                \ell(T) \times r(T) - o(T) = z(T)h(T).
            \end{equation*}
            Denote $\boldsymbol{\ell}(T) = (\ell_1(T),\dots,\ell_n(T))$, and similarly $\boldsymbol{r}(T)$, $\boldsymbol{o}(T)$.
            \pause

            Now, consider the following \textbf{linear PCP for QAP}:
            \begin{algoen}
                \item $\mathcal{P}$ commits to an extended witness $\mathbf{z}$ and coefficients $\mathbf{h} = (h_1,\dots,h_n)$ of $h(T)$.\pause
                \item $\mathcal{V}$ samples $\gamma \leftarrowS \mathbb{F}$ and sends query $\boldsymbol{\gamma} = (\gamma,\gamma^2,\dots,\gamma^n)$ to $\mathcal{P}$.\pause
                \item $\mathcal{P}$ reveals the following values:
                \begin{align*}
                    \pi_1 &\gets \langle \mathbf{z}, \boldsymbol{\ell}(\gamma) \rangle, \quad &\pi_2 &\gets \langle \mathbf{z}, \boldsymbol{r}(\gamma) \rangle, \\
                    \pi_3 &\gets \langle \mathbf{z}, \boldsymbol{o}(\gamma) \rangle, \quad &\pi_4 &\gets z(\gamma) \cdot \langle \mathbf{h}, \boldsymbol{\gamma} \rangle.
                \end{align*}
                \vspace{-10pt}
                \pause
                \item $\mathcal{V}$ checks whether $\pi_1\pi_2 - \pi_3 = \pi_4$.
            \end{algoen}
        \end{example}
    \end{frame}

    \begin{frame}{Why is it Sound?}
        \begin{questionblock}{Question}
            Why is it safe to use such a check? (assuming proper commitments).
        \end{questionblock}
        \pause

        The polynomials $\ell(T)$, $r(T)$ and $o(T)$ are interpolated polynomials using
        $\left| \Circ \right|$ (number of gates) points, so:
        \vspace{-5pt}
        \begin{equation*}
            \deg(\ell) \le \left| \Circ \right|, \quad
            \deg(r) \le \left| \Circ \right|, \quad
            \deg(o) \le \left| \Circ \right|
            \vspace{-5pt}
        \end{equation*}
        \pause
        Thus, we can estimate the degree of polynomial $m(T) = \ell(T)r(T) - o(T)$.
        \vspace{-5pt}
        \begin{equation*}
            \deg(m) \leq \max\{\deg(\ell) + \deg(r), \deg(o)\} \leq 2 \left| \Circ \right|
            \vspace{-5pt}
        \end{equation*}
        \pause
        If an adversary $\mathcal{A}$ does not know a valid witness $\mathbf{z}$, he can compute a
        polynomial $(\widetilde{m}(T), \widetilde{h}(T)) \gets \mathcal{A}(\cdot)$ that satisfies a verifier
        $\mathcal{V}$ with probability:
        \vspace{-5pt}
        \begin{equation*}
            \Pr_{s \leftarrowS \mathbb{F}}[\widetilde{m}(s) = z(s)\widetilde{h}(s)] \leq \frac{2 \left| \Circ \right|}{\#\mathbb{F}}
        \end{equation*}

        If $\#\mathbb{F}$ is large enough, $2|\Circ|/\#\mathbb{F}$ is \textit{negligible}.
    \end{frame}

    \section{Proof Of Exponent}

    \begin{frame}{Encrypted Verification}
        Let us try to prove that we know some polynomial $p(T)$ that can be divided by $t(T)$ without a
        remainder.
        \pause

        Consider the polynomial: $p(T) = T^2 - 5T + 2$. Additionally, we will need a cyclic group $\mathbb{G}$ with a generator $g \in \mathbb{G}$. We also define the \textit{encryption} operation as follows:
        \begin{equation*}
            \mathsf{Enc}: \mathbb{F} \to \mathbb{G}, \quad \mathsf{Enc}(x) := g^x
        \end{equation*}
        \pause

        Essentially, $\mathsf{Enc}(p(\tau))$ is the \textbf{KZG Commitment}. Let us see the encryption of $p(\tau)$ for our example:
        \begin{equation*}
            \mathsf{Enc}(p(\tau)) = g^{p(\tau)} = g^{\left( \tau^2 - 5\tau + 2 \right)}  = \left( g^{\tau^2} \right)^1 \cdot {\left( g^{\tau^1} \right)}^{-5} \cdot \left(g^{\tau^0}\right)^2
        \end{equation*}
        \pause

        \begin{remarkblock}{Note}
            KZG Commitment requires only encrypted powers of $\tau$: $\{g^{\tau^i}\}_{i \in [d]}$.
        \end{remarkblock}
    \end{frame}

    \begin{frame}{Encrypted Verification}
        \textbf{Verifier}:
        \vspace{-5pt}
        \begin{itemize}[label=\ding{51}]
            \item \vspace{-3pt} Picks a random value $\tau \leftarrowS \mathbb{F}$. \pause
            \item \vspace{-3pt} Calculates the public parameters $\{g^{\tau^i}\}_{i \in [d]}$. \pause
            \item \vspace{-3pt} Calculates $t(\tau)$. \pause
            \item \vspace{-3pt} Outputs prover parameters $\{g^{\tau^i}\}_{i \in [d]}$. \pause
        \end{itemize}
        \textbf{Prover}:
        \vspace{-5pt}
        \begin{itemize}[label=\ding{51}]
            \item \vspace{-3pt} Calculates $h(T) = \frac{p(T)}{t(T)}$. \pause
            \item \vspace{-3pt} Using $\{g^{\tau^i}\}_{i \in [d]}$ calculates $g^{p(\tau)}$ and $g^{h(\tau)}$. \pause
            \item \vspace{-3pt} Provides encrypted polynomials $g^{p(\tau)}$ and $g^{h(\tau)}$ to the verifier. \pause
        \end{itemize}
        \textbf{Verifier}:
        \vspace{-5pt}
        \begin{itemize}[label=\ding{51}]
            \item \vspace{-3pt} Checks whether $g^{p(\tau)} = \left( g^{h(\tau)} \right)^{t(\tau)}$.
        \end{itemize}
    \end{frame}

    \begin{frame}{That doesn't work...}
        \textbf{Verifier}:
        \vspace{-5pt}
        \begin{itemize}[label=\ding{51}]
            \item \vspace{-3pt} Picks a random value $\tau \leftarrowS \mathbb{F}$. \pause
            \item \vspace{-3pt} Calculates the public parameters $\{g^{\tau^i}\}_{i \in [d]}$. \pause
            \item \vspace{-3pt} Calculates $t(\tau)$. \pause
            \item \vspace{-3pt} Outputs prover parameters $\{g^{\tau^i}\}_{i \in [d]}$. \pause
        \end{itemize}
        \textbf{Adversary}:
        \vspace{-5pt}
        \begin{itemize}[label=\ding{51}]
            \item \vspace{-3pt} Picks a random value $r \leftarrowS \mathbb{F}$, calculates $g^r$. \pause
            \item \vspace{-3pt} Calculates $g^{t(\tau)}$. \pause
            \item \vspace{-3pt} Calculates $g^{\widetilde{p}(\tau)} = {\left( g^{t(\tau)} \right)}^r$. \pause
        \end{itemize}
        \textbf{Verifier}:
        \vspace{-5pt}
        \begin{itemize}[label=\ding{51}]
            \item \vspace{-3pt} Checks whether $g^{\widetilde{p}(\tau)} = \left( g^{r} \right)^{t(\tau)}$.
        \end{itemize}
    \end{frame}

    \begin{frame}{Proof Of Exponent}
        \textbf{Verifier}:
        \vspace{-5pt}
        \begin{itemize}[label=\ding{51}]
            \item \vspace{-3pt} Picks random values $\tau \leftarrowS \mathbb{F}$, $\alpha \leftarrowS \mathbb{F}$. \pause
            \item \vspace{-3pt} Calculates the public parameters $\{g^{\tau^i}\}_{i \in [d]}$ and $\{g^{\alpha\tau^i}\}_{i \in [d]}$. \pause
            \item \vspace{-3pt} Calculates $t(\tau)$. \pause
        \end{itemize}
        \textbf{Prover}:
        \vspace{-5pt}
        \begin{itemize}[label=\ding{51}]
            \item \vspace{-3pt} Calculates $h(T) = \frac{p(T)}{t(T)}$. \pause
            \item \vspace{-3pt} Using $\{g^{\tau^i}\}_{i \in [d]}$ calculates $g^{p(\tau)}$, $g^{h(\tau)}$. \pause
            \item \vspace{-3pt} Using $\{g^{\alpha\tau^i}\}_{i \in [d]}$ calculates $g^{p'(\tau)} = g^{\alpha p(\tau)}$. \pause
            \item \vspace{-3pt} Provides encrypted polynomials to the verifier. \pause
        \end{itemize}
        \textbf{Verifier}:
        \vspace{-5pt}
        \begin{itemize}[label=\ding{51}]
            \item \vspace{-3pt} Checks whether $g^{p(\tau)} = \left( g^{h(\tau)} \right)^{t(\tau)}$. \pause
            \item \vspace{-3pt} Checks whether $g^{p'(\tau)} = \left( g^{p(\tau)} \right)^{\alpha} = g^{\alpha p(\tau)}$.
        \end{itemize}
    \end{frame}

    \makeclosingframe
\end{document}
